\begin{abstract}
Alignment via instruction tuning is widely believed to reduce the diversity
of language model outputs, imposing an ``alignment tax'' on stylistic and
creative capabilities. We test this assumption directly by measuring
stylistic impersonation ability across seven literary and functional styles
in a paired base/aligned model setup (\qwenbase and \qweninst). We then
apply logit-space distribution arithmetic---linearly interpolating base and
aligned model logits at each decoding step---to investigate whether base
model distributions survive alignment and remain continuously navigable.
Our results challenge the naive alignment tax narrative: the aligned model
\emph{significantly outperforms} the base model on style fidelity
($7.26$ vs.\ $3.74$, $p < 0.0001$), and is the only condition that
reliably differentiates between target styles (ANOVA $F\!=\!6.32$,
$p\!=\!0.002$). Distribution arithmetic produces smooth, monotonic
interpolation between base and aligned behavior, confirming that base
distributions are neither destroyed nor discontinuously altered by
alignment. We reframe the alignment--diversity tradeoff: base models
possess diverse latent distributions but lack the controllability to
access them on demand; alignment adds an addressing mechanism
(instruction following) that unlocks this access. The relevant axis is
not preservation versus loss, but controllability versus distributional
freedom.
\end{abstract}
